\section*{Devoir maison}

\begin{Exercise}[title={Boucles imbriquées}]

\begin{lstlisting}[style=pythonstyle, language=python]
    x = int(input("Saisissez un nombre entier :"))
    acc = 0
    for i in range(x):
       for j in range(i):
            acc = (i-j)**2
    print(f"Le total vaut {acc}")
\end{lstlisting}

\Question Donnez les variables et leurs types.
\Question Complétez le tableau d'avancement pour une saisie de x = 2.\\
\begin{center}
    \large
\begin{tabular}{c|c|c|c|c|}
\textbf{Ligne} & \textbf{x} & \textbf{i} & \textbf{j} & \textbf{acc}  \\
\hline
 & &  &  & \\
 & &  &  & \\
 & &  &  & \\
  & &  &  & \\
   & &  &  & \\
    & &  &  & \\
     & &  &  & \\
\end{tabular}
\end{center}

\begin{center}
\scalebox{0.7}{

{\VideoBox{Introduction a la programmation/DM flash/BLOB/f-srting video/qrcode.png}{Introduction a la programmation/DM flash/BLOB/f-srting video/thumbnail.png}{Pour comprendre la ligne 6 et les f-strings en python.~\url{https://www.youtube.com/watch?v=RgzX4qrLmPM}}}
}
\end{center}


\Question Qu'affichera le code pour :
\begin{itemize}
    \item x = 5 :\\
    \item x = 12 :\\
    \item x = 120 :\\
    \item x = -1\\
\end{itemize}
\end{Exercise}
\begin{Exercise}[title={La factorielle}]
\Question On souhaite maintenant calculer la factorielle d'un nombre x. Par exemple la factorielle de 5 noté 5! = 5*4*3*2*1. Écrire au dos de la copie le code permettant de réaliser le calcul pour n'importe quel nombre x donné par l'utilisateur.
\end{Exercise}